\begin{frame}
\frametitle{Obiettivo dei test} 
\begin{itemize}
\item Valutare la capacità di RITA nel rilevare i Covert Channel ICMP
\item Analizzare l'impatto che dei parametri hanno sulla gravità del canale: 
\begin{itemize}
\item Quantità di dati inviati
\item Campi usati nei messai ICMP 
\item Uso di un delay tra i pacchetti
\item Tempo di attesa prima di un nuovo invio 
\item Modifica del mittente nel pacchetto
\end{itemize} 
%\item Identificare quali implementazioni sono più difficili da rilevare
\end{itemize} 
\end{frame} 

%\begin{frame}
%\frametitle{Test 1: Singolo invio}
%\textbf{Setup}
%\begin{itemize}
%\item Invio dati una sola volta
%\item Blocchi da 1KB
%\item Delay variabile (1–15s)
%\end{itemize}
%\vspace{1em}
%\textbf{Risultato}
%\begin{itemize}
%\item RITA non rileva il covert channel
%\end{itemize}
%\vspace{1em}
%\textbf{Interpretazione}
%\begin{itemize}
%\item Attività troppo sporadica
%\item Nessun pattern identificabile
%\item Assenza di pattern temporali
%\item Non classificata come beaconing
%\end{itemize}
%\end{frame}

%\begin{frame}
%\frametitle{Test 2: Invii multipli dei dati} 
%5\begin{itemize}
%\item I dati vengono inviati più volte (3 volte)
%\item Uso dello stesso covert channel
%\item Variazione dei parametri
%\end{itemize}
%\vspace{1em}
%Si vogliono rilevare i pattern che influenzano la gravità del 
%Covert Channel. %Generare pattern rilevabili 
%\end{frame}

\begin{frame}
\frametitle{Risultati}
Tempi di attesa: 
\begin{itemize}
\item Lunghi tempi di attesa riducono la gravità
%\item La quantità di dati che si inviano in un secondo è minore
\end{itemize} 
Invio sequenziale dei pacchetti 
\begin{itemize} 
\item Trasmissione non rilevata da RITA 
\item Genera rumore sul canale di trasmissione. %Traffico anomalo sulla rete
\end{itemize}
\vspace{1em} 
Quantità di dati 
\begin{itemize}
\item Una quantità ridotta (100KB) spesso non viene rilevata
%\item Possibile sotto soglia di detection 
\item Una quantità elevata (1MB) spesso viene rilevata 
\end{itemize}
%\vspace{1em} 
%Payload variabile: 
%\begin{itemize} 
%\item Non influenza molto la rilevabilità
%%\item Non migliora la stealthiness
%%\item Beacon score fino al 79\%
%%\item Possibile peggioramento rilevabilità
%\end{itemize} 
\vspace{1em} 
Modifica nei pacchetti del mittente 
\begin{itemize}
\item Riduce la gravità del canale 
\item Può aumentare il beacon score (in base a come viene implementato)
\end{itemize}
\end{frame}

%\begin{frame}
%\frametitle{Risultati}
%Effetto del tempo di attesa: 
%\begin{itemize}
%\item 15 minuti → possibile riduzione rilevabilità
%\item 30 minuti → aumento beacon score
%\item 1 ora → comportamento intermedio
%\end{itemize} 
%Test senza delay 
%\begin{itemize}
%\item Invio sequenziale dei dati
%\item RITA non rileva il covert channel
%\end{itemize}
%\vspace{1em}
%\textbf{Limite}
%\begin{itemize}
%\item Traffico anomalo sulla rete
%\end{itemize}
%Test con 100KB
%\begin{itemize}
%\item Dati ridotti → spesso non rilevati
%\item Possibile sotto soglia di detection
%\end{itemize}
%Payload variabile: 
%\begin{itemize}
%\item Non migliora la stealthiness
%\item Beacon score fino al 79\%
%\item Possibile peggioramento rilevabilità
%\end{itemize}
%Spoofing del mittente
%\begin{itemize}
%\item Riduce la gravità (Low vs Medium)
%\item Può aumentare il beacon score
%\item Dipende da come viene implementato
%\end{itemize}
%\end{frame}


\begin{frame}
\frametitle{Considerazioni}
\begin{itemize}
\item Il pattern temporale è critico 
\item Il contenuto è stato meno rilevante rispetto al comportamento
%\item Il delay influisce sulla rilevabilità
%\item Lo modifica dei mittenti deve essere controllata. 
%%riduce la gravita ma aumenta il beaconing se implementato male. 
%\item Il traffico troppo veloce può evitare detection ma crea anomalie
\end{itemize} 
\vspace{1em} 
Riduce la rilevabilità: 
\begin{itemize}
\item Modifica nei pacchetti del mittente
\item Tempi di attesa lunghi tra gli invii 
\item Schemi di invio irregolari
%\item Invii non frequenti e sporadici
%\item Riduzione gravità
%\item Riduzione beacon score
\end{itemize}
Aumenta la rilevabilità: 
\begin{itemize}
\item Invii frequenti o ravvicinati
\item Schemi di invio regolari
%\item Dimensione payload variabile in modo evidente
%\item Aumento gravità
%\item Aumento beacon score
\end{itemize}
\end{frame} 

%\begin{frame}
%\frametitle{Conclusione}
%\begin{itemize}
%\item La detection dipende dal pattern temporale
%\item Il contenuto è meno rilevante del comportamento
%\item Tecniche per ridurre detection possono introdurre anomalie di rete
%\end{itemize}
%\end{frame}

%\begin{frame}
%\frametitle{Insight chiave}
%\begin{itemize}
%\item Il pattern temporale è più importante del contenuto
%\item Il delay influisce direttamente sulla rilevabilità
%\item Lo spoofing riduce la gravità ma può aumentare il beaconing
%\end{itemize}
%\end{frame}

%\begin{frame}
%\frametitle{Risultati principali}
%\textbf{Riduce la rilevabilità}
%\begin{itemize}
%\item invii sporadici
%\item spoofing mittente
%\item tempi lunghi tra invii
%\end{itemize}
%\vspace{1em}
%\textbf{Aumenta la rilevabilità}
%\begin{itemize}
%\item invii frequenti
%\item pattern regolari
%\item variazioni evidenti nel payload
%\end{itemize}
%\end{frame}



%\begin{frame} [fragile]
%\frametitle{Test e Risultati}  
%\scriptsize 
%\begin{longtable}{|p{0.45\textwidth}|p{0.45\textwidth}|}
%    \multicolumn{2}{c}{Primo Test} \\ \hline 
%    \textbf{Descrizione} & \textbf{Risultato} \\ \hline 
%    Si inviano i dati una signola volta & RITA non ha rilevato il Covert Channel \\ \hline 
%    \multicolumn{2}{c}{} \\[2ex] 
%    \multicolumn{2}{c}{Secondo Test} \\ \hline 
%    \textbf{Descrizione} & \textbf{Risultato} \\ \hline 
%    Si inviano i dati più volte (in particolare 3) al variare di alcuni parametri\par 
%    \begin{itemize}
%    \item quantità di dati inviati
%    \item campi utilizzati nel messaggio
%    \item delay durante l'invio dei dati 
%    \item attesa prima di un nuovo invio
%    \item mittente reale o modificato
%    \end{itemize}
%    Usato lo stesso Covert Channel per ogni test
%    & 
%    Riducono la gravita della comunicazione ed il tasso di beacon associato: 
%    \begin{itemize}
%    \item modifica del mittente 
%    \item tempi di attesa lunghi prima dell'invio successivo  
%    \end{itemize} 
%    \vspace{1ex} 
%    Aumentano la gravità ed il beacon score: 
%    \begin{itemize}
%    \item l'invio costante, ravvicinato e frequente di messaggi 
%    \item una variazione scostante della dimensione del campo Data
%    \end{itemize} 
%    \\ \hline 
%    \caption{Mappatura degli eventi ai dati}  
%\end{longtable} 
%\end{frame} 

